\chapter{Inteligência artificial de baixo consumo para o bem-estar estudantil}

Este capítulo retoma, em português, a nota técnica
\emph{Low-Power Artificial Intelligence for Monitoring Student Well-Being and
Burnout Risk in Low-Connectivity Contexts in Brazil}. O texto examina de que
maneira a proposta poderia funcionar nas redes públicas brasileiras, levando em
conta as condições de infraestrutura, as responsabilidades dos envolvidos e os
limites da intervenção.

\section{Cenário atual}

As escolas públicas brasileiras enfrentam níveis crescentes de ansiedade,
estresse crônico e exaustão acadêmica entre os estudantes. Esses fenômenos
prejudicam o engajamento, o desempenho e a permanência escolar. No Brasil, o problema é agravado pelo estigma associado à saúde mental. A
análise suplementar do PISA usada na nota técnica indica que aproximadamente
70\% dos jovens evitam admitir que estão estressados. O dado converge com a
enquete do UNICEF, na qual medo, insegurança e falta de informação aparecem
entre os motivos para não procurar ajuda \cite{unicef2022ajuda}.

Também há pouca capacidade institucional para perceber e responder a essas
situações. A nota técnica registra que somente 15,7\% das escolas públicas
brasileiras teriam psicólogo presente na unidade. Na prática, a identificação
tende a depender do julgamento informal de professores que já acumulam muitas
responsabilidades. A Lei nº 13.935/2019 prevê serviços de psicologia e serviço
social nas redes públicas, mas a existência da norma não significa que esses
profissionais estejam disponíveis em todas as escolas e territórios
\cite{brasil13935}.

A conectividade precária não causa o sofrimento psicológico, mas amplia a
dificuldade de resposta. No cenário analisado pela nota técnica, 63,7\% das
escolas da região Norte estavam conectadas à internet. Plataformas educacionais
concebidas para uso contínuo da nuvem tornam-se pouco viáveis onde a conexão é
intermitente. Os dados da TIC Educação 2024 reforçam que o acesso nominal não
elimina limitações de velocidade, dispositivos e cobertura dentro das escolas
\cite{cgi2025}.

Os dados mais recentes da TIC Educação ajudam a qualificar esse diagnóstico.
Embora o acesso nominal à internet tenha avançado, persistem problemas de
velocidade, cobertura interna, manutenção e disponibilidade de dispositivos
\cite{cgi2025}. Portanto, a política não deve equiparar ``escola conectada'' a
``escola capaz de manter um serviço em nuvem''. A qualidade e a continuidade do
acesso são relevantes para a escolha tecnológica.

\section{Solução proposta}

Enfrentar essa lacuna exige uma ferramenta projetada para as condições que as
escolas efetivamente possuem, e não para a infraestrutura presumida pelos
fornecedores. Propõe-se um aplicativo móvel \emph{offline-first} que aplique um
instrumento psicossocial validado: a DASS-21, ou Escala de Depressão, Ansiedade
e Estresse. Para adolescentes brasileiros, a referência adequada é a EDAE-A,
versão adaptada e validada por Patias e colaboradores \cite{patias2016}.

Em vez de enviar as respostas a um servidor em nuvem, o aplicativo realiza a
pontuação localmente e utiliza um modelo leve de classificação do tipo TinyML.
Essa forma de IA de baixo consumo segue princípios de baixa conectividade e
hardware mínimo presentes na abordagem AIED Unplugged \cite{uema2026}. O
estudante responde ao questionário em um dispositivo compartilhado da escola; o
sistema organiza as respostas em faixas auxiliares de atenção; e a equipe recebe
um alerta protegido.

Para a gestão, a informação deve ser agregada e anonimizada, como no exemplo:
``cinco estudantes desta turma apresentaram indicação de atenção moderada ou
elevada''. Nomes e respostas individuais ficam fora dos painéis gerenciais. O
acesso individualizado deve ser limitado a profissional treinado e autorizado,
quando necessário para realizar a escuta e dar continuidade ao cuidado.

O processamento local procura tornar a triagem mais regular sem transferir a
decisão ao aplicativo. O alerta apenas organiza a demanda. A interpretação das
respostas, o acolhimento e a definição do encaminhamento continuam sob
responsabilidade dos profissionais.

A agregação e a confidencialidade são requisitos centrais. Mesmo um sistema
bem-intencionado pode produzir efeito contrário ao pretendido caso o estudante
acredite que será rotulado, punido ou exposto. Nessa situação, ele pode ocultar
o sofrimento ou fornecer respostas socialmente desejáveis. A proteção não é
apenas uma exigência jurídica; ela determina a confiança e a qualidade dos dados
obtidos.

\begin{figure}[htbp]
  \centering
  \resizebox{\textwidth}{!}{%
  \begin{tikzpicture}[
    node distance=7mm and 10mm,
    every node/.style={font=\scriptsize,align=center},
    caixa/.style={draw,rounded corners=3pt,minimum height=10mm,
      text width=48mm,inner sep=3pt},
    fluxo/.style={-{Latex[length=2mm]},gray!75,semithick}
  ]
    \node[font=\small\bfseries] (titulo)
      {IA de baixo consumo --- Brasil (TinyML offline)};
    \node[caixa,fill=violet!10,below left=7mm and 2mm of titulo] (educacao)
      {Ministério da Educação\\Diretrizes, LGPD e recursos};
    \node[caixa,fill=violet!10,below right=7mm and 2mm of titulo] (saude)
      {Ministério da Saúde\\Diretrizes de saúde mental};
    \node[caixa,fill=green!10,below=13mm of titulo] (gestores)
      {Gestores e diretores escolares\\Implantam alertas e lideram a cultura de cuidado};
    \node[caixa,fill=orange!12,below left=9mm and 2mm of gestores] (professores)
      {Professores\\Interpretam alertas e intervêm};
    \node[caixa,fill=green!12,below right=9mm and 2mm of gestores] (estudantes)
      {Estudantes\\Dados da DASS-21 e classificação de risco};
    \node[caixa,fill=red!8,below=14mm of gestores] (aplicativo)
      {Aplicativo TinyML offline\\Processamento local em dispositivos de baixo custo};
    \node[caixa,fill=gray!10,below left=9mm and 2mm of aplicativo] (familias)
      {Famílias\\Consentimento e adoção legítima};
    \node[caixa,fill=gray!10,below right=9mm and 2mm of aplicativo] (tecnologia)
      {Setor tecnológico e TIC\\Desenvolvimento da ferramenta offline};

    \draw[fluxo] (educacao) -- (gestores);
    \draw[fluxo] (saude) -- (gestores);
    \draw[fluxo] (gestores) -- (professores);
    \draw[fluxo] (gestores) -- (estudantes);
    \draw[fluxo] (professores) -- (aplicativo);
    \draw[fluxo] (estudantes) -- (aplicativo);
    \draw[fluxo] (aplicativo) -- (familias);
    \draw[fluxo] (aplicativo) -- (tecnologia);
  \end{tikzpicture}
  }
  \caption{Atores e relações operacionais da proposta de IA de baixo consumo}
  \label{fig:atores-tinyml}
  \par\smallskip
  \small Fonte: tradução e adaptação de \emph{catedra\_texto.pdf}.
\end{figure}

A Figura \ref{fig:atores-tinyml} mostra a articulação prevista. Educação e saúde
definem diretrizes, recursos e regras; gestores escolares conduzem a implantação
e promovem uma cultura de cuidado; professores interpretam alertas e ajudam a
acionar o fluxo; estudantes fornecem os dados da DASS-21; famílias participam da
legitimação e do cuidado; e parceiros tecnológicos apoiam o desenvolvimento da
solução offline.

\section{Cenário político}

O Brasil já deu um passo legislativo ao instituir a Política Nacional de Atenção
Psicossocial nas Comunidades Escolares por meio da Lei nº 14.819/2024. A norma
reconhece a saúde mental como componente da qualidade educacional e determina a
articulação entre educação, saúde e assistência social \cite{brasil14819}. A
lacuna encontra-se agora na capacidade de implementar o que a política prevê.

A execução, porém, depende de coordenação entre setores e de capacidade local
para acolher os estudantes. A nota técnica aponta um déficit acumulado de
R\$ 117,71 bilhões no chamado Orçamento do Conhecimento desde 2014. Esse dado é
usado no documento para situar as restrições de financiamento que afetam as
redes de ensino. A lei atribui a governança aos Grupos de Trabalho
Intersetoriais do Programa Saúde na Escola e exige a participação da saúde e da
comunidade escolar \cite{brasil14819}. O programa proposto deve partir dessa
estrutura, em vez de criar um fluxo paralelo.

A nota técnica também registra redução de 2,3\% nas matrículas da educação
básica em 2025. Como esse movimento pode decorrer de mudanças demográficas e
administrativas, ele não deve ser interpretado automaticamente como evasão. Os
resultados do programa precisam ser comunicados com a mesma cautela,
distinguindo variação ao longo do tempo, associação e efeito causal.

\section{Atores envolvidos}

Os ministérios e as secretarias de Educação e Saúde ocupam posição central:
definem diretrizes, distribuem recursos e estabelecem regras para que a
ferramenta cumpra a LGPD e as políticas nacionais de saúde mental escolar. Seu
interesse é elevado, pois o tema possui relevância social e potencial retorno
público. A liderança deve ser compartilhada com a assistência social e com as
instâncias territoriais do Programa Saúde na Escola.

Diretores e professores participam da implantação cotidiana e da interpretação
dos alertas. A adesão desses profissionais dependerá de o sistema ser percebido
como apoio, e não como nova forma de cobrança. Isso exige formação para respostas
emocionais, informacionais e instrumentais, além de reorganização das tarefas.
Se apenas acrescentar trabalho, a tecnologia tende a ser abandonada ou utilizada
de maneira superficial.

Os estudantes são o grupo mais afetado. Seu engajamento depende, sobretudo, da
confiança de que os dados não serão utilizados para rotular, discriminar ou
punir. Pais e responsáveis também influenciam a legitimidade do programa. A
comunicação de privacidade deve explicar, em linguagem simples, o que é coletado,
quem pode acessar e quais providências podem ocorrer.

Parceiros tecnológicos, universidades e organizações da sociedade civil podem
fornecer capacidade para construir e adaptar o sistema às condições locais. A
nota técnica menciona organizações especializadas em ferramentas
\emph{offline-first} para o Sul Global. Essas parcerias devem estar subordinadas
à governança pública, com padrões abertos, portabilidade, transparência e
proibição de exploração comercial dos dados estudantis.

\section{Desafios e lacunas}

A proposta enfrenta problemas de naturezas diferentes. Boa parte das soluções
de IA pressupõe uma infraestrutura que muitas escolas não possuem. Ao mesmo
tempo, o subfinanciamento e a falta de profissionais limitam a resposta ao que a
triagem revela. Há ainda o estigma, que pode levar estudantes a esconder o que
sentem mesmo quando existe um canal de escuta.

Nenhuma dessas lacunas é resolvida apenas com um modelo mais preciso. O desafio
tecnológico pede arquitetura local e manutenção simples; o institucional exige
governança, formação e articulação da rede; o sociocultural requer comunicação,
confidencialidade e participação estudantil. É por esse motivo que as
recomendações técnicas são apresentadas em conjunto com medidas institucionais e
de comunicação.

\section{Comparação das alternativas de política}

A nota técnica compara o aplicativo offline com duas alternativas: contratação
universal de profissionais e plataforma dependente de alta conectividade. A
comparação não pretende afirmar que uma alternativa elimina as demais. A
contratação de psicólogos e assistentes sociais continua necessária; o
aplicativo atua como apoio à organização da demanda.

\begin{longtable}{p{3.0cm}p{3.3cm}p{3.3cm}p{3.3cm}}
\caption{Comparação das alternativas apresentadas no artigo}\label{tab:alternativas-artigo}\\
\toprule
\textbf{Alternativa} & \textbf{Adequação técnica} & \textbf{Apoio político} & \textbf{Viabilidade organizacional} \\
\midrule
\endfirsthead
\toprule
\textbf{Alternativa} & \textbf{Adequação técnica} & \textbf{Apoio político} & \textbf{Viabilidade organizacional} \\
\midrule
\endhead
Aplicativo de classificação \emph{offline-first}
& Alta: a solução de baixo consumo responde à limitação de conectividade.
& Média-alta: há interesse público no bem-estar, mas preocupações com privacidade podem gerar oposição.
& Média: requer parceria especializada em TinyML e gestão da carga de trabalho. \\
\addlinespace
Contratação universal de profissionais
& Alta: oferece apoio humano direto e possui fundamento científico.
& Média: custo recorrente e sustentabilidade orçamentária podem gerar resistência.
& Baixa em áreas remotas: há dificuldade de contratar e reter profissionais qualificados. \\
\addlinespace
Plataforma de alta conectividade
& Alta em ambientes com infraestrutura e dados adequados.
& Média-baixa: requer investimento elevado e pode criar dependência tecnológica.
& Baixa em escolas sem internet estável, especialmente em territórios historicamente excluídos. \\
\bottomrule
\end{longtable}

\section{Recomendações}

O artigo recomenda uma política apoiada em IA de baixo consumo,
\emph{offline-first} e TinyML para o monitoramento do bem-estar em contextos de
baixa conectividade. O processamento local em dispositivos de baixo custo amplia
o alcance potencial da política. Sua implementação é organizada em quatro ações
coordenadas.

\subsection{Implementação prática}

A implementação começaria pelo desenvolvimento do aplicativo móvel e por um
piloto de alcance limitado. O sistema coleta dados de bem-estar por meio da DASS-21/EDAE-A, um
questionário de 21 itens que avalia sintomas recentes de depressão, ansiedade e
estresse. Seu uso é de rastreio e monitoramento, sem função diagnóstica. Modelos
leves processam localmente os dados e geram faixas auxiliares de atenção em um
sistema de alertas agregados e anonimizados. A ação da escola continua sujeita
ao julgamento profissional.

O artigo sugere iniciar com aproximadamente cem escolas das regiões Norte e
Nordeste, em articulação com iniciativas de AIED Unplugged e TinyML. A meta
inicial proposta é redução de 15\% nos escores da escala entre a linha de base e
o momento posterior à intervenção. Essa porcentagem é uma hipótese mensurável,
não um resultado garantido. Os dados do piloto serviriam para validar o
desempenho, corrigir protocolos, reduzir riscos e fundamentar eventual expansão
nacional.

\subsection{Institucionalização das rotinas de resposta}

A resposta humana precisa fazer parte da rotina escolar antes que a tecnologia
seja ampliada. Módulos de formação continuada devem preparar professores e equipes
para compreender alertas e oferecer apoio emocional, informacional e
instrumental. Também é necessário impedir que a ferramenta se torne mais um
item na rotina já sobrecarregada: reorganizar tarefas, retirar obrigações
administrativas sem relação com o cuidado e criar incentivos institucionais
influencia a percepção do sistema como apoio ou vigilância.

\subsection{Monitoramento orientado à equidade e à transparência}

A confiança no programa também depende de regras de dados que possam ser
compreendidas e fiscalizadas. Isso inclui armazenamento local cifrado, acesso individual
restrito a profissionais treinados, relatórios agregados e anonimizados para
gestores e proibição clara do uso punitivo ou discriminatório. Essas medidas não
são detalhes posteriores da implantação: constituem uma condição para a
participação dos estudantes e das famílias.

\subsection{Tradução da pesquisa}

Os resultados da pesquisa ainda precisam ser adaptados à rotina das escolas.
Adaptar a DASS-21 ao uso cotidiano da escola, produzir orientações claras para
profissionais não especialistas e acompanhar indicadores quantitativos e
qualitativos permite que a ferramenta seja aprimorada continuamente. Entre os
indicadores quantitativos estão os escores da escala e a permanência escolar;
entre os qualitativos, a percepção da comunidade sobre confiança, acolhimento e
utilidade.

A IA de baixo consumo não resolve a falta de apoio psicossocial nas escolas.
Ela pode ajudar a organizar a procura por atendimento e a identificar situações
que exigem atenção, sobretudo onde a conexão é instável. O resultado, porém,
depende do que acontece depois do alerta: se não houver acolhimento e
encaminhamento, o sistema apenas registrará uma necessidade que a rede continua
sem condições de atender.
